\chapter*{Agradecimientos}
Es complicado ser un punto en un hiperplano lleno de temas por aprender y habilidades por descubrir. Cuando entré a la carrera de matemáticas, no me imaginaba el mundo de cosas que me esperaban. Desde mi primer semestre me enamoré por áreas muy teóricas y me enemisté con la computación. El haber conocido el club Pu++ y con ello la olimpiada de programación ICPC, significó para mí el renacer de un interés por aprender computación y en la actualidad mi rama de trabajo. Agradezco profundamente a los fundadores de Pu++ y todos los nombres que de forma altruista decidieron dedicar su tiempo a enseñar a más personas como instructores, moldearon una gran parte de quién soy hoy y mi camino.
\\\\
Agradezco a mis padres, Lucy Téllez y Román Melo y a mis hermanos Xavier Melo y JuanPablo Hernández, así como a mi tía Juana Téllez porque aunque nunca comprendieron qué estudiaba en mi carrera no les impidió apoyarme y darme su más sincera ayuda. Sin ellos nada de la carrera hubiera podido ser posible.
\\\\
A Manuel Alcántara por aceptar con tanta emoción un trabajo de docencia no tan ortodoxo, por guiarme a lo largo de este trabajo.
\\\\
Agradezco a mis amigos de la facultad Alex, Beto, Vin, Héctor, Rodrigo, Mario, Ana Pau, Daniel, Paco porque la carrera es una montaña rusa de emociones en las que muchas veces uno se siente solo y estuvieron cuando más lo necesitaba.
\\\\
Agradezco también a los profesores que más me marcaron a lo largo de la carrera con sus enseñanzas en clases, así como ideas filósoficas. A Martha Takane y Judith Alvarado por ayudarme a no sentirme sola en una carrera dominada por hombres. A Angel Cuahutemoc por ayudarme a darme cuenta mi gusto por la teoría de números. A Mauricio Gabriel Medina por haberme despertado el amor por el álgebra moderna. A César Cruz por haberme adoptado académicamente aún cuando me decidiera por seguir otro camino lejano a la investigación.
\\\\
Agradezco a Simba, que ya no se encuentra conmigo pero estuvo conmigo muchas noches de la carrera mientras hacia una demostración que no me salía dandome apoyo. 
\\\\
A Nico por apoyarme esta recta final, y acompañarme tantas tardes trabajando y editando este trabajo, emocionada por la vida que nos espera.
\\\\
A todos los concursantes de ICPC que pusieron su tiempo en resolver los problemas y darme retroalimentación, gracias, no hay nada peor que un problema nunca testeado!

